% Options for packages loaded elsewhere
\PassOptionsToPackage{unicode}{hyperref}
\PassOptionsToPackage{hyphens}{url}
%
\documentclass[
  ignorenonframetext,
  aspectratio=43,
]{beamer}
\usepackage{pgfpages}
\setbeamertemplate{caption}[numbered]
\setbeamertemplate{caption label separator}{: }
\setbeamercolor{caption name}{fg=normal text.fg}
\beamertemplatenavigationsymbolshorizontal
% Prevent slide breaks in the middle of a paragraph
\widowpenalties 1 10000
\raggedbottom
\setbeamertemplate{part page}{
  \centering
  \begin{beamercolorbox}[sep=16pt,center]{part title}
    \usebeamerfont{part title}\insertpart\par
  \end{beamercolorbox}
}
\setbeamertemplate{section page}{
  \centering
  \begin{beamercolorbox}[sep=12pt,center]{part title}
    \usebeamerfont{section title}\insertsection\par
  \end{beamercolorbox}
}
\setbeamertemplate{subsection page}{
  \centering
  \begin{beamercolorbox}[sep=8pt,center]{part title}
    \usebeamerfont{subsection title}\insertsubsection\par
  \end{beamercolorbox}
}
\AtBeginPart{
  \frame{\partpage}
}
\AtBeginSection{
  \ifbibliography
  \else
    \frame{\sectionpage}
  \fi
}
\AtBeginSubsection{
  \frame{\subsectionpage}
}

\usepackage{amsmath,amssymb}
\usepackage{lmodern}
\usepackage{iftex}
\ifPDFTeX
  \usepackage[T1]{fontenc}
  \usepackage[utf8]{inputenc}
  \usepackage{textcomp} % provide euro and other symbols
\else % if luatex or xetex
  \usepackage{unicode-math}
  \defaultfontfeatures{Scale=MatchLowercase}
  \defaultfontfeatures[\rmfamily]{Ligatures=TeX,Scale=1}
\fi
\usetheme[]{default}
\usecolortheme{default}
% Use upquote if available, for straight quotes in verbatim environments
\IfFileExists{upquote.sty}{\usepackage{upquote}}{}
\IfFileExists{microtype.sty}{% use microtype if available
  \usepackage[]{microtype}
  \UseMicrotypeSet[protrusion]{basicmath} % disable protrusion for tt fonts
}{}
\makeatletter
\@ifundefined{KOMAClassName}{% if non-KOMA class
  \IfFileExists{parskip.sty}{%
    \usepackage{parskip}
  }{% else
    \setlength{\parindent}{0pt}
    \setlength{\parskip}{6pt plus 2pt minus 1pt}}
}{% if KOMA class
  \KOMAoptions{parskip=half}}
\makeatother
\usepackage{xcolor}
\newif\ifbibliography
\setlength{\emergencystretch}{3em} % prevent overfull lines
\setcounter{secnumdepth}{-\maxdimen} % remove section numbering


\providecommand{\tightlist}{%
  \setlength{\itemsep}{0pt}\setlength{\parskip}{0pt}}\usepackage{longtable,booktabs,array}
\usepackage{calc} % for calculating minipage widths
\usepackage{caption}
% Make caption package work with longtable
\makeatletter
\def\fnum@table{\tablename~\thetable}
\makeatother
\usepackage{graphicx}
\makeatletter
\def\maxwidth{\ifdim\Gin@nat@width>\linewidth\linewidth\else\Gin@nat@width\fi}
\def\maxheight{\ifdim\Gin@nat@height>\textheight\textheight\else\Gin@nat@height\fi}
\makeatother
% Scale images if necessary, so that they will not overflow the page
% margins by default, and it is still possible to overwrite the defaults
% using explicit options in \includegraphics[width, height, ...]{}
\setkeys{Gin}{width=\maxwidth,height=\maxheight,keepaspectratio}
% Set default figure placement to htbp
\makeatletter
\def\fps@figure{htbp}
\makeatother

\usepackage{ragged2e}
\usepackage{etoolbox}

\apptocmd{\frame}{}{\justifying}{} 

\setbeamertemplate{enumerate items}[default]
\setbeamertemplate{itemize items}[circle]
\setbeamersize{text margin left=2.25em, text margin right=2.25em}

\let\OLDitemize\itemize
\renewcommand\itemize{
    \OLDitemize\addtolength{\itemsep}{3pt plus 1pt minus 1pt}
}
\let\OLDenumerate\enumerate
\renewcommand\enumerate{
    \OLDenumerate\addtolength{\itemsep}{3pt plus 1pt minus 1pt}
}

\usepackage{orcidlink}
\usepackage{amsmath}
\DeclareMathOperator*{\argmax}{argmax}

\makeatletter
\def\@listii{\leftmargin\leftmarginii
              \topsep    1ex
              \parsep    0\p@   \@plus\p@
              \itemsep   \parsep}
\makeatother

\usepackage{lipsum}

\let\olditem\item
\renewcommand\item{\olditem\justifying}

\definecolor{blue}{RGB}{0,114,178}
\definecolor{red}{RGB}{213,94,0}
\definecolor{yellow}{RGB}{240,228,66}
\definecolor{green}{RGB}{0,158,115}

\definecolor{MyBackground}{RGB}{255,253,218}

\setbeamercolor{section in toc}{fg=blue}
\setbeamercolor{subsection in toc}{fg=red}

\setbeamercolor{frametitle}{fg=blue}
\setbeamercolor{title}{fg=blue}

\setbeamertemplate{footline}[frame number]
\setbeamertemplate{navigation symbols}{} 
\setbeamertemplate{itemize items}{-}

\setbeamercolor{itemize item}{fg=blue}
\setbeamercolor{itemize subitem}{fg=blue}
\setbeamercolor{enumerate item}{fg=blue}
\setbeamercolor{enumerate subitem}{fg=blue}
\setbeamercolor{button}{bg=MyBackground,fg=blue}

\newcounter{daggerfootnote}
\newcommand*{\daggerfootnote}[1]{%
    \setcounter{daggerfootnote}{\value{footnote}}%
    \renewcommand*{\thefootnote}{\fnsymbol{footnote}}%
    \footnote[2]{#1}%
    \setcounter{footnote}{\value{daggerfootnote}}%
    \renewcommand*{\thefootnote}{\arabic{footnote}}%
}

\AtBeginSection{
    \begingroup
    \setbeamercolor{background canvas}{bg=yellow}
    \begin{frame}
        \begin{centering}
            \usebeamerfont{section title}\huge\color{black}\insertsection\par
        \end{centering}
    \end{frame}
    \endgroup
}
\makeatletter
\makeatother
\makeatletter
\makeatother
\makeatletter
\@ifpackageloaded{caption}{}{\usepackage{caption}}
\AtBeginDocument{%
\ifdefined\contentsname
  \renewcommand*\contentsname{Table of contents}
\else
  \newcommand\contentsname{Table of contents}
\fi
\ifdefined\listfigurename
  \renewcommand*\listfigurename{List of Figures}
\else
  \newcommand\listfigurename{List of Figures}
\fi
\ifdefined\listtablename
  \renewcommand*\listtablename{List of Tables}
\else
  \newcommand\listtablename{List of Tables}
\fi
\ifdefined\figurename
  \renewcommand*\figurename{Figure}
\else
  \newcommand\figurename{Figure}
\fi
\ifdefined\tablename
  \renewcommand*\tablename{Table}
\else
  \newcommand\tablename{Table}
\fi
}
\@ifpackageloaded{float}{}{\usepackage{float}}
\floatstyle{ruled}
\@ifundefined{c@chapter}{\newfloat{codelisting}{h}{lop}}{\newfloat{codelisting}{h}{lop}[chapter]}
\floatname{codelisting}{Listing}
\newcommand*\listoflistings{\listof{codelisting}{List of Listings}}
\makeatother
\makeatletter
\@ifpackageloaded{caption}{}{\usepackage{caption}}
\@ifpackageloaded{subcaption}{}{\usepackage{subcaption}}
\makeatother
\makeatletter
\@ifpackageloaded{tcolorbox}{}{\usepackage[many]{tcolorbox}}
\makeatother
\makeatletter
\@ifundefined{shadecolor}{\definecolor{shadecolor}{rgb}{.97, .97, .97}}
\makeatother
\makeatletter
\makeatother
\ifLuaTeX
  \usepackage{selnolig}  % disable illegal ligatures
\fi
\IfFileExists{bookmark.sty}{\usepackage{bookmark}}{\usepackage{hyperref}}
\IfFileExists{xurl.sty}{\usepackage{xurl}}{} % add URL line breaks if available
\urlstyle{same} % disable monospaced font for URLs
\hypersetup{
  pdftitle={Seminar 9},
  pdfauthor={Gabriel E. Cabrera Guzmán},
  hidelinks,
  pdfcreator={LaTeX via pandoc}}

\title[Seminar 9]{Seminar
9\daggerfootnote{\texttt{gabriel.cabreraguzman@postgrad.manchester.ac.uk}}}
\subtitle{Statistics 7 (Forecasting) Solutions}
\author[Gabriel E. Cabrera Guzmán]{
    {Gabriel E. Cabrera
Guzmán\textsuperscript{\normalsize\orcidlink{0000-0000-0000-0000}}} \medskip
}
\date{December 13, 2025}
\institute[]{
    The University of Manchester\\
Alliance Manchester Business School \bigskip
}

\definecolor{quarto-callout-note-color}{RGB}{0,114,178}
\definecolor{quarto-callout-note-color-frame}{RGB}{0,114,178}
\definecolor{quarto-callout-warning-color}{RGB}{213,94,0}
\definecolor{quarto-callout-warning-color-frame}{RGB}{213,94,0}
\begin{document}
\frame{\titlepage}
\ifdefined\Shaded\renewenvironment{Shaded}{\begin{tcolorbox}[interior hidden, borderline west={3pt}{0pt}{shadecolor}, boxrule=0pt, sharp corners, breakable, enhanced, frame hidden]}{\end{tcolorbox}}\fi

\begin{frame}{Exercise 1}
\protect\hypertarget{exercise-1}{}
Exponential Forecasting with \(\alpha = 0.2\): \vspace{0.15cm}

\begin{table}[h!]
\centering
\begin{tabular}{c c c c}
\toprule
\textbf{Month} & $\mathbf{x_t}$ & $\mathbf{F_t}$ & $\mathbf{e_t = x_t - F_t}$ \\
\midrule
1  & 23 &        &        \\
2  & 39 & 40.00  &        \\
3  & 36 &        &        \\
4  & 26 &        &        \\
5  & 36 &        &        \\
6  & 12 &        &        \\
7  & 26 &        &        \\
8  & 30 &        &        \\
9  & 42 &        &        \\
10 & 35 &        &        \\
11 & 36 &        &        \\
12 & 28 &        &        \\
\bottomrule
\end{tabular}
\end{table}
\end{frame}

\begin{frame}{Exercise 1 (Cont'd)}
\protect\hypertarget{exercise-1-contd}{}
\begin{enumerate}
\item
  Complete the above table for exponential forecasting assuming
  \(\alpha = 0.2\). \vspace{0.15cm}

  \pause

  \begin{table}[h!]
   \centering
   \begin{tabular}{c c c c}
   \toprule
   \textbf{Month} & $\mathbf{x_t}$ & $\mathbf{F_t}$ & $\mathbf{e_t = x_t - F_t}$ \\
   \midrule
   1  & 23 &        &         \\
   2  & 39 & 40.00  &  -1.00  \\
   3  & 36 & 39.80  &  -3.80  \\
   4  & 26 & 39.04  & -13.04  \\
   5  & 36 & 36.43  &  -0.43  \\
   6  & 12 & 36.35  & -24.35  \\
   7  & 26 & 31.48  &  -5.48  \\
   8  & 30 & 30.38  &  -0.38  \\
   9  & 42 & 30.30  &  11.70  \\
   10 & 35 & 32.64  &   2.36  \\
   11 & 36 & 33.12  &   2.88  \\
   12 & 28 & 33.69  &  -5.69  \\
   \bottomrule
   \end{tabular}
   \end{table}
\end{enumerate}
\end{frame}

\begin{frame}{Exercise 1 (Cont'd)}
\protect\hypertarget{exercise-1-contd-1}{}
\begin{enumerate}
\setcounter{enumi}{1}
\item
  State assumption(s) required for exponential forecasting.

  \pause

  The time series should be absent of trend and seasonal effect
\end{enumerate}
\end{frame}

\begin{frame}{Exercise 2}
\protect\hypertarget{exercise-2}{}
\begin{table}[h!]
\centering
\begin{tabular}{c c c c}
\toprule
\textbf{Month} & $\mathbf{x_t}$ & $\mathbf{F_t}$ & $\mathbf{e_t = x_t - F_t}$ \\
\midrule
1 & 45 &        &        \\
2 & 62 & 67.00  &        \\
3 & 73 &        &        \\
4 & 87 &        &        \\
5 & 86 &        &        \\
6 & 95 &        &        \\
\bottomrule
\end{tabular}
\end{table}
\end{frame}

\begin{frame}{Exercise 2 (Cont'd)}
\protect\hypertarget{exercise-2-contd}{}
\begin{enumerate}
\item
  Complete the above table for exponential forecasting assuming
  \(\alpha = 0.2\) and \(\alpha = 0.4\).

  \pause

  When \(\alpha = 0.2\), we have: \vspace{0.25cm}

  \pause

  \begin{table}[h!]
   \centering
   \begin{tabular}{c c c c}
   \toprule
   Month & $\mathbf{x_t}$ & $\mathbf{F_t}$ & $\mathbf{e_t = x_t - F_t}$ \\
   \midrule
   1 & 45 &        &        \\
   2 & 62 & 67.00  & -5.00  \\
   3 & 73 & 66.00  & 7.00   \\
   4 & 87 & 67.40  & 19.60  \\
   5 & 86 & 71.32  & 14.68  \\
   6 & 95 & 74.26  & 20.74  \\
   \bottomrule
   \end{tabular}
   \end{table}
\end{enumerate}
\end{frame}

\begin{frame}{Exercise 2 (Cont'd)}
\protect\hypertarget{exercise-2-contd-1}{}
\pause

When \(\alpha = 0.4\), we have: \vspace{0.25cm}

\pause

\begin{table}[h!]
\centering
\begin{tabular}{c c c c}
\toprule
Month & $\mathbf{x_t}$ & $\mathbf{F_t}$ & $\mathbf{e_t = x_t - F_t}$ \\
\midrule
1 & 45 &        &        \\
2 & 62 & 67.00  & -5.00  \\
3 & 73 & 65.00  & 8.00   \\
4 & 87 & 68.20  & 18.80  \\
5 & 86 & 75.72  & 10.28  \\
6 & 95 & 79.83  & 15.17  \\
\bottomrule
\end{tabular}
\end{table}
\end{frame}

\begin{frame}{Exercise 2 (Cont'd)}
\protect\hypertarget{exercise-2-contd-2}{}
\begin{enumerate}
\setcounter{enumi}{1}
\item
  Compute the Mean Squared Error (MSE) assuming \(\alpha = 0.2\) and
  \(\alpha = 0.4\).

  \pause

  MSE when \(\alpha = 0.2\): \vspace{-0.5cm}

  \pause

  \[
   \frac{(-5)^2 + 7^2 + (19.6)^2 + (14.68)^2 + (20.74)^2}{5} = \frac{4492.22}{5} = 898.44
   \]

  \pause

  MSE when \(\alpha = 0.4\): \vspace{-0.5cm}

  \pause

  \[
   \frac{(-5)^2 + 8^2 + (18.8)^2 + (10.28)^2 + (15.17)^2}{5} = \frac{3277.33}{5} = 655.47
   \]
\end{enumerate}
\end{frame}

\begin{frame}{Exercise 2 (Cont'd)}
\protect\hypertarget{exercise-2-contd-3}{}
\begin{enumerate}
\setcounter{enumi}{2}
\item
  Compute the Mean Absolute Error (MAE) assuming \(\alpha = 0.2\) and
  \(\alpha = 0.4\).

  \pause

  MAE when \(\alpha = 0.2\): \vspace{-0.5cm}

  \pause

  \[
   \frac{|{-5}| + |7| + |19.6| + |14.68| + |20.74|}{5} = \frac{67.02}{5} = 13.40
   \]

  \pause

  MAE when \(\alpha = 0.4\): \vspace{-0.5cm}

  \pause

  \[
   \frac{|{-5}| + |8| + |18.8| + |10.28| + |15.17|}{5} = \frac{57.25}{5} = 11.45
   \]
\end{enumerate}
\end{frame}

\begin{frame}{Exercise 2 (Cont'd)}
\protect\hypertarget{exercise-2-contd-4}{}
\begin{enumerate}
\setcounter{enumi}{3}
\item
  Based on (2) and (3), suggest which \(\alpha\) is more suitable for
  exponential forecasting.

  \pause

  Both MSE and MAE are lower when \(\alpha = 0.4\). Therefore,
  \(\alpha = 0.4\) is a better choice for exponential forecasting.
\end{enumerate}
\end{frame}

\begin{frame}{Exercise 3}
\protect\hypertarget{exercise-3}{}
\begin{table}[h!]
\centering
\begin{tabular}{c c c c c c c c}
\toprule
\textbf{Period} & $\mathbf{x_t}$ &
\textbf{Period} & $\mathbf{x_t}$ &
\textbf{Period} & $\mathbf{x_t}$ &
\textbf{Period} & $\mathbf{x_t}$ \\
\midrule
1  & 22   & 8  & 20   & 15 & 22   & 22 & 22 \\
2  & 18   & 9  & 25   & 16 & 24   & 23 & 28 \\
3  & 15   & 10 & 21   & 17 & 29   & 24 & 31 \\
4  & 19   & 11 & 20   & 18 & 27   &     &    \\
5  & 23   & 12 & 20   & 19 & 21   &     &    \\
6  & 21   & 13 & 27   & 20 & 26   &     &    \\
7  & 16   & 14 & 25   & 21 & 29   &     &    \\
\bottomrule
\end{tabular}
\end{table}
\end{frame}

\begin{frame}{Exercise 3 (cont'd)}
\protect\hypertarget{exercise-3-contd}{}
\pause

\begin{enumerate}
\item
  Compute 3-month simple moving average for the above data.
\item
  Compute 4-month simple moving average for the above data.
\item
  Compute centered 4-month simple moving average for the above data.
\end{enumerate}
\end{frame}



\end{document}
